\usepackage{listings}
\usepackage{float}
\usepackage{fancyhdr}
\usepackage[pdftex]{graphicx}
\usepackage{amsmath}
\usepackage{nicefrac}
\usepackage{amsthm}
\usepackage{amssymb}
\usepackage{listing}
\usepackage{enumitem}
\usepackage{apptools}
\usepackage{chngcntr}

%%% Allow math equation to cross pages

%
%
%

%
\usepackage{mathtools}
% \usepackage{bbm}
\usepackage{wrapfig}
%
\usepackage{color}
%
%
%
%
%
%
%
%
%
\usepackage[ruled, vlined, linesnumbered]{algorithm2e}
\DeclareMathOperator*{\argmin}{argmin}
\DeclareMathOperator*{\argmax}{argmax}

\usepackage[capitalise,nameinlink,sort]{cleveref}
\Crefname{equation}{Eq.}{Eqs.}
\Crefname{assumption}{Assumption}{Assumptions}
\Crefname{condition}{Condition}{Conditions}
\Crefname{problem}{Linear Program}{Linear Programs}
\Crefname{question}{Questions}{Questions}

\usepackage{quoting}

\usepackage[parfill]{parskip}

\usepackage{tikz}

\tikzstyle{round}=[thick,draw=black,circle]

\usepackage{thm-restate}

\theoremstyle{plain}
\newtheorem{theorem}{\textbf{Theorem}}
\newtheorem{lemma}{\textbf{Lemma}}
\newtheorem{corollary}[lemma]{Corollary}
\newtheorem{proposition}[lemma]{Proposition}
\newtheorem{claim}[lemma]{Claim}
\newtheorem{observation}[lemma]{Observation}
\newtheorem{conjecture}[lemma]{Conjecture}
\newtheorem*{lemma-a}{\textbf{Lemma}}
\newtheorem*{theorem-a}{\textbf{Theorem}}
\newtheorem{question}{Question}
\newtheorem{example}{Example}
\newtheorem{definition}{Definition}

\theoremstyle{definition}
\newtheorem{assumption}{Assumption}
\newtheorem{fact}{Fact}
\newtheorem{remark}{Remark}
\hypersetup{colorlinks=true,citecolor=blue,linkcolor=red}
\setcitestyle{citesep={;}}
\DeclareMathOperator*{\dprime}{\prime \prime}
